\documentclass[UTF8,fontset=fandol,zihao=5]{ctexart}

\usepackage[a4paper,margin=2cm]{geometry}
\usepackage{amsmath,amssymb,mathtools}
\usepackage{upgreek}
\usepackage[nolimits]{cmupint}
\usepackage{booktabs,tabularx,array}
\usepackage{xcolor}
\usepackage[most]{tcolorbox}
\usepackage{enumitem}
\usepackage{indentfirst}
\usepackage{hyperref}
\usepackage{float}

\usepackage{tikz}
\usepackage{diagbox}

\newcolumntype{W}[1]{%
		>{\hsize=#1\hsize\linewidth=\hsize
		\centering\arraybackslash}X}

\usetikzlibrary{
		positioning,
		calc,
		decorations.pathreplacing,
		arrows.meta
}

% relationbox 中知识关系图的统一样式
\tikzset{
		relnode/.style={
		font=\small,
		align=center
		},
		relline/.style={
		thick
		},
		relarrow/.style={
		-{Stealth},
		thick
		}
}

\DeclareTextFontCommand{\textbf}{\heiti\color{main}}

\definecolor{main}{HTML}{D0645A}
\definecolor{mainlight}{HTML}{FFF0F0}
\definecolor{method}{HTML}{D09A95}
\definecolor{methodlight}{HTML}{fff0f0}
\definecolor{black}{HTML}{000000}

\hypersetup{colorlinks=true,linkcolor=main,urlcolor=main}
\setlength{\parindent}{2em}
\linespread{1.5}
\setlength{\abovedisplayskip}{5pt}
\setlength{\belowdisplayskip}{5pt}
\setlength{\abovedisplayshortskip}{3pt}
\setlength{\belowdisplayshortskip}{3pt}
\setlist[itemize]{leftmargin=2em,itemsep=0.1em,topsep=0.2em}
\setlist[enumerate]{leftmargin=2.3em,itemsep=0.1em,topsep=0.2em}

\newfontfamily\headingfont{TeX Gyre Heros}

\ctexset{
		section={format=\centering\heiti\headingfont\zihao{3}\color{main},
				name={第,章},number=\chinese{section},aftername=\quad,
				beforeskip=0.9em,afterskip=0.5em},
		subsection={format=\heiti\headingfont\zihao{5}\color{main},
				name={,、},number=\chinese{subsection},aftername=,
				beforeskip=0.6em,afterskip=0.22em},
		subsubsection={format=\heiti\headingfont\zihao{5}\color{main},
				number={（\chinese{subsubsection}）},aftername=,
				beforeskip=0.55em,afterskip=0.15em}
}

\newtcolorbox{relationbox}{
	enhanced,breakable,colback=mainlight,colframe=main,
	boxrule=0.65pt,arc=1.4mm,left=2mm,right=2mm,top=1.2mm,bottom=1.2mm
}
\newtcolorbox{methodbox}[1][]{
	enhanced,breakable,colback=methodlight,colframe=method,
	boxrule=0.65pt,arc=1.4mm,left=2mm,right=2mm,top=1.2mm,bottom=1.2mm,
	title=#1,fonttitle=\kaishu,fontupper=\kaishu,
	before upper={\setlength{\parindent}{2em}}
}

\newcommand{\jj}{\,\mathrm{j}}
\newcommand{\ee}{\,\mathrm{e}}
\newcommand{\CC}{\mathbb{C}}
\newcommand{\RR}{\mathbb{R}}
\newcommand{\QQ}{\mathbb{Q}}
\newcommand{\dd}{\,\mathrm{d}}
\newcommand{\ppi}{\,\uppi\,}
\newcommand{\FF}{\mathcal{F}}
\newcommand{\LL}{\mathcal{L}}
\newcommand{\Sa}{\operatorname{Sa}}
\newcommand{\RRe}{\operatorname{Re}}
\newcommand{\IIm}{\operatorname{Im}}
\newcommand{\Res}{\operatorname{Res}}
\let\bm\boldsymbol

\newfontfamily\russianfont{DejaVu Sans}

\begin{document}

\begin{center}
	{\heiti\headingfont\zihao{1}\color{main} 矢量分析与场论笔记}
\end{center}

\section{场量}

\subsection{矢量微分算子}

\subsubsection{梯度$\nabla \varphi$}

\subsubsection{旋度$\nabla\times\bm A$}

\subsubsection{散度$\nabla\cdot\bm A$}

\subsubsection{Laplace算子$\nabla^2\bm A$}

\subsection{矢量积分特征}

\subsubsection{环量$\displaystyle\varGamma=\oint_C \bm A\cdot\dd\bm l$}

\subsubsection{通量$\displaystyle\varPhi=\oiint_S\bm A\cdot\dd\bm S$}

\section{场论定理}

\subsection{广义Stokes定理}

\subsection{Holmholtz定理}

定义Hamilton算子（Nabla算子）
\[
\nabla=\left(\frac{\partial}{\partial x_1},
\frac{\partial}{\partial x_2},\ldots,
\frac{\partial}{\partial x_n}\right),
\]

Laplace 算子
\[
\Delta=\nabla^2=\sum_{i=1}^{n}\frac{\partial^2}{\partial x_i^2}.
\]

\subsection{环量与通量}

\textbf{（1）环量}

作闭合曲线 $C$，定义环量
\[
\Gamma=\oint_C\bm A\cdot\dd\bm s.
\]
保证法向量方向不变的前提下，曲线 $C$ 以任意方式缩小为点 $M$ 时，若极限
\[
\lim_{\Delta S\to0}\frac{\Delta\Gamma}{\Delta S}
\]
存在，则定义环量密度
\[
\frac{\dd\Gamma}{\dd S}\bm n.
\]

\textbf{（2）通量}

作闭合曲面 $S$，定义通量
\[
\Phi=\oiint_S\bm A\cdot\dd\bm S.
\]
若 $\Phi=0$，说明曲面内部无源或汇，否则有源或汇。

曲面 $S$ 以任意方式缩小为点 $M$ 时，若极限
\[
\lim_{\Delta V\to0}\frac{\Delta\Phi}{\Delta V}
\]
存在，则定义通量密度（即散度）
\[
\operatorname{div}\bm A=\nabla\cdot\bm A
=\frac{\dd\Phi}{\dd V}.
\]

\textbf{2. 梯度、旋度与散度}

\textbf{（1）梯度}

梯度是向量值，指示数量场在某点上升最快的速度和方向，用 $\nabla f$ 表示。

线性：
\[
\nabla(C_1u+C_2v)=C_1\nabla u+C_2\nabla v.
\]
运算性质：
\[
\nabla(uv)=u\nabla v+v\nabla u,
\qquad
\nabla\left(\frac uv\right)=\frac{v\nabla u-u\nabla v}{v^2}.
\]
复合法则：
\[
\nabla f(u)=f'(u)\nabla u.
\]

\textbf{（2）旋度}

旋度是向量值，指示向量场在某点的旋转强度与方向，用 $\nabla\times\bm A$ 表示。数量上等于该点环量密度的最大值，方向即环量密度最大方向。

线性：
\[
\nabla\times(C_1\bm A\pm C_2\bm B)
=C_1\nabla\times\bm A\pm C_2\nabla\times\bm B.
\]
运算性质：
\[
\nabla\times(u\bm A)=u\nabla\times\bm A+\nabla u\times\bm A.
\]

\textbf{（3）散度}

散度是数量值，指示向量场在某点发散或汇聚的强度，用 $\nabla\cdot\bm A$ 表示。散度即通量密度。

线性：
\[
\nabla\cdot(C_1\bm A\pm C_2\bm B)
=C_1\nabla\cdot\bm A\pm C_2\nabla\cdot\bm B.
\]
运算性质：
\[
\nabla\cdot(u\bm A)=u\nabla\cdot\bm A+\nabla u\cdot\bm A.
\]

\textbf{（4）其他计算公式}

梯度场无旋：
\[
\nabla\times(\nabla u)=\bm0.
\]
旋度场无源：
\[
\nabla\cdot(\nabla\times\bm A)=0.
\]
向量积散度公式：
\[
\nabla\cdot(\bm A\times\bm B)
=\bm B\cdot(\nabla\times\bm A)-\bm A\cdot(\nabla\times\bm B).
\]
梯度场散度公式：
\[
\nabla\cdot(\nabla u)=\nabla^2u.
\]
旋度场旋度公式：
\[
\nabla\times(\nabla\times\bm A)
=\nabla(\nabla\cdot\bm A)-\nabla^2\bm A.
\]

\textbf{3. 特殊向量场}

\textbf{（1）无旋场}

单连通区域内，无旋场、有势场、保守场等价。设 $\bm A\in C^1(G)$，$G\subset\RR^3$：
\[
\begin{array}{ll}
\text{无旋场：}&\nabla\times\bm A=\bm0,\\
\text{有势场：}&\exists u,\ \text{s.t. }\bm A=\nabla u,\\
\text{保守场：}&\displaystyle\oint_C\bm A\cdot\dd\bm s=0.
\end{array}
\]

积分与路径无关等价于保守场、旋度为 $0$，即等价于 $\nabla\times\bm A=\bm0$。

\textbf{（2）无源场}

单连通区域内，无源场、向量有势场、保守场等价：
\[
\begin{array}{ll}
\text{无源场：}&\nabla\cdot\bm A=0,\\
\text{向量有势场：}&\exists\bm B,\ \text{s.t. }\bm A=\nabla\times\bm B,\\
\text{保守场：}&\displaystyle\oiint_S\bm A\cdot\dd\bm S=0.
\end{array}
\]

\textbf{（3）调和场}

调和场势函数满足 Laplace 方程
\[
\nabla^2u=0,
\]
即
\[
\frac{\partial^2u}{\partial x^2}
+\frac{\partial^2u}{\partial y^2}
+\frac{\partial^2u}{\partial z^2}=0.
\]

\subsection{定理}

\textbf{0. Newton--Leibniz 公式}
\[
\int_a^b f'(x)\dd x=f(b)-f(a).
\]

\textbf{1. Green 公式}

设平面有界闭域 $\sigma$ 由分段光滑的简单闭曲线围成，$\sigma$ 的边界曲线记为 $C$，函数 $P,Q\in C^1(\sigma)$，则
\[
\iint_{\sigma}\left(\frac{\partial Q}{\partial x}-\frac{\partial P}{\partial y}\right)\dd x\dd y
=\oint_{C^+}P\dd x+Q\dd y.
\]

\textbf{2. Stokes 公式}
\[
\iint_S\left(\frac{\partial R}{\partial y}-\frac{\partial Q}{\partial z}\right)\dd y\wedge\dd z
+\left(\frac{\partial P}{\partial z}-\frac{\partial R}{\partial x}\right)\dd z\wedge\dd x
+\left(\frac{\partial Q}{\partial x}-\frac{\partial P}{\partial y}\right)\dd x\wedge\dd y
=\oint_{C^+}P\dd x+Q\dd y+R\dd z.
\]

几何意义为旋度对面积的积分等于原函数在外曲线上的积分，即向量形式
\[
\iint_S(\nabla\times\bm A)\cdot\dd\bm S
=\oint_{C^+}\bm A\cdot\dd\bm s.
\]

\textbf{3. Gauss 公式}

设空间有界闭域 $V$ 由分段光滑的简单闭曲面 $S$ 围成，向量函数 $\bm A=(P,Q,R)\in C^1(V)$，则
\[
\iiint_V\left(\frac{\partial P}{\partial x}
+\frac{\partial Q}{\partial y}
+\frac{\partial R}{\partial z}\right)\dd V
=\oiint_{S^+}P\dd y\wedge\dd z
+Q\dd z\wedge\dd x
+R\dd x\wedge\dd y.
\]

几何意义为散度对体积的积分等于原函数在外表面上的积分，即向量形式
\[
\iiint_V\nabla\cdot\bm A\dd V
=\oiint_{S^+}\bm A\cdot\dd\bm S.
\]

上三式可以统一于如下公式，受限于知识水平不作进一步介绍。

\textbf{4. 广义 Stokes 公式}
\[
\int_{\Omega}\dd\omega=\int_{\partial\Omega}\omega.
\]

\begin{methodbox}[广义 Stokes 公式求第二型曲线曲面积分]
广义 Stokes 公式是解题一把好手。使用广义 Stokes 公式解题时应注意以下三点：

\begin{enumerate}
	\item 原函数处处有定义，尤其是原点；
	\item 路径为正；
	\item 原函数可导。
\end{enumerate}

此外，还要注意正负与方向：曲线左侧为积分区域时为正，否则为负；曲面外侧为正，内侧为负。常用方法为右手定则。

常用方法有二：

\textbf{填补法}：将不完整的曲线或曲面补为闭合曲线或曲面，通过
\[
\int_{\Omega}\bm A\cdot\dd\bm\omega
=\int_{\Sigma^+}\bm A\cdot\dd\bm\omega
+\int_{\Sigma^-}\bm A\cdot\dd\bm\omega,
\]
即用闭合曲线或曲面积分减去添加曲线或曲面积分，得到原始路径积分。该方法通常有两种特征：

\begin{enumerate}
	\item 图形便于填补，如半圆、扇形弧，柱体侧面、锥体侧面、各种旋转体侧面；
	\item 图形补充后闭合部分散度为 $0$，即 $\nabla\cdot\bm A=0$。
\end{enumerate}

\textbf{挖空法}：用于定义域内存在间断点的情况，以原点居多。此时需用到路径变形定理。

二维表述：在无旋场中，内部包含奇点的同伦同向闭合曲线积分值相同。为将原函数表达式中形如
\[
\frac{x^2}{a^2}+\frac{y^2}{b^2}
\]
的分母部分消掉，作小椭圆
\[
\begin{cases}
x=a\varepsilon\cos\theta,\\
y=b\varepsilon\sin\theta,
\end{cases}
\]
$\varepsilon\to0$ 时，同向椭圆路径上的积分值即原积分值。常见函数如
\[
\bm A=\left(-\frac{y}{x^2+y^2},\frac{x}{x^2+y^2}\right).
\]

三维表述：在无源场内，内部包含奇点的同伦同向闭合曲面积分值相同。为将原函数表达式中形如
\[
\frac{x^2}{a^2}+\frac{y^2}{b^2}+\frac{z^2}{c^2}
\]
的分母部分消掉，作小椭球体
\[
\begin{cases}
x=a\varepsilon\sin\varphi\cos\theta,\\
y=b\varepsilon\sin\varphi\sin\theta,\\
z=c\varepsilon\cos\varphi,
\end{cases}
\]
$\varepsilon\to0$ 时，同向椭球闭合曲面上的积分值即原积分值。
\end{methodbox}

\begin{methodbox}[由向量函数求势函数]
主要用偏积分法。即分别对 $x$、对 $y$ 积分，将二者合并得到势函数。注意要将不定积分中的 $C$ 改为 $C(x)$、$C(y)$。
\end{methodbox}

\end{document}